\section{Unggah File (File Upload) dan Validasi}

Fitur \textbf{file upload} memungkinkan pengguna mengirimkan file dari browser ke server melalui form HTML. Ini digunakan dalam banyak aplikasi: mengunggah foto profil, lampiran dokumen, bukti pembayaran, atau impor data CSV. Form upload memerlukan dua syarat khusus: metode \code{POST} dan atribut \code{enctype="multipart/form-data"} pada elemen \code{<form>}. Elemen input menggunakan tipe \code{file}: \code{<input type="file" name="berkas">}. Tanpa \code{enctype} yang benar, file tidak akan terkirim ke server \cite{phpManual}.

Di sisi server, data file yang diunggah tersedia di superglobal \code{\$\_FILES}. Array ini berisi informasi penting: \code{name} (nama file asli), \code{type} (tipe MIME yang dilaporkan browser), \code{size} (ukuran dalam byte), \code{tmp\_name} (lokasi sementara di server), dan \code{error} (kode error). File yang berhasil diunggah disimpan sementara di direktori temporary server dan harus dipindahkan ke lokasi permanen menggunakan fungsi \code{move\_uploaded\_file()} sebelum skrip berakhir \cite{phpRightWay}.

Validasi file upload sangat krusial untuk keamanan. Beberapa langkah validasi yang wajib dilakukan: (1) periksa kode error: \code{\$\_FILES['berkas']['error'] === UPLOAD\_ERR\_OK}; (2) batasi ukuran file (misalnya maks. 2 MB); (3) validasi tipe file menggunakan \textit{whitelist} ekstensi yang diizinkan; (4) jangan percaya \code{\$\_FILES['type']} dari client --- validasi MIME di server menggunakan \code{finfo\_file()} atau \code{mime\_content\_type()}; (5) buat nama file baru (misalnya dengan \code{uniqid()}) untuk mencegah overwrite atau path traversal. Kegagalan validasi upload dapat membuat server rentan terhadap serangan \textit{file inclusion} atau \textit{remote code execution} \cite{owaspXss}.

\begin{figure}[ht]
\centering
\begin{tikzpicture}[
  box/.style={draw, rounded corners, minimum width=2.8cm, minimum height=0.9cm, align=center, font=\footnotesize, fill=#1},
  decision/.style={draw, diamond, aspect=2, minimum width=2.2cm, align=center, font=\footnotesize, fill=yellow!15},
  arrow/.style={-Stealth, thick},
  node distance=1cm
]
  \node[box=blue!15] (start) {User pilih\\file \& submit};
  \node[box=green!15, below=of start] (receive) {PHP terima\\\code{\$\_FILES}};
  \node[decision, below=of receive] (err) {Error\\= 0?};
  \node[box=red!15, right=2.5cm of err] (reject1) {Tolak:\\error upload};
  \node[decision, below=1.2cm of err] (size) {Ukuran\\$\leq$ maks?};
  \node[box=red!15, right=2.5cm of size] (reject2) {Tolak:\\terlalu besar};
  \node[decision, below=1.2cm of size] (type) {Tipe\\valid?};
  \node[box=red!15, right=2.5cm of type] (reject3) {Tolak:\\tipe tidak\\diizinkan};
  \node[box=purple!15, below=1.2cm of type] (move) {\code{move\_uploaded\_file()}\\ke folder aman};
  \node[box=green!20, below=of move] (done) {Upload berhasil};

  \draw[arrow] (start) -- (receive);
  \draw[arrow] (receive) -- (err);
  \draw[arrow] (err) -- node[above, font=\scriptsize]{Tidak} (reject1);
  \draw[arrow] (err) -- node[left, font=\scriptsize]{Ya} (size);
  \draw[arrow] (size) -- node[above, font=\scriptsize]{Tidak} (reject2);
  \draw[arrow] (size) -- node[left, font=\scriptsize]{Ya} (type);
  \draw[arrow] (type) -- node[above, font=\scriptsize]{Tidak} (reject3);
  \draw[arrow] (type) -- node[left, font=\scriptsize]{Ya} (move);
  \draw[arrow] (move) -- (done);
\end{tikzpicture}
\caption{Flowchart validasi file upload}
\end{figure}

\begin{webcode}[language=PHP, caption={Upload file dengan validasi lengkap}]
<?php
if ($_SERVER['REQUEST_METHOD'] === 'POST' && isset($_FILES['berkas'])) {
  $file     = $_FILES['berkas'];
  $maxSize  = 2 * 1024 * 1024; // 2 MB
  $allowed  = ['jpg', 'jpeg', 'png', 'pdf'];
  $uploadDir = __DIR__ . '/uploads/';

  // 1. Cek error
  if ($file['error'] !== UPLOAD_ERR_OK) {
    die("Error upload: kode " . $file['error']);
  }

  // 2. Cek ukuran
  if ($file['size'] > $maxSize) {
    die("File terlalu besar (maks. 2 MB).");
  }

  // 3. Cek ekstensi (whitelist)
  $ext = strtolower(pathinfo($file['name'], PATHINFO_EXTENSION));
  if (!in_array($ext, $allowed)) {
    die("Tipe file tidak diizinkan. Hanya: " . implode(', ', $allowed));
  }

  // 4. Validasi MIME di server
  $finfo = finfo_open(FILEINFO_MIME_TYPE);
  $mime  = finfo_file($finfo, $file['tmp_name']);
  finfo_close($finfo);
  $allowedMime = ['image/jpeg', 'image/png', 'application/pdf'];
  if (!in_array($mime, $allowedMime)) {
    die("Tipe MIME tidak valid: {$mime}");
  }

  // 5. Buat nama unik dan pindahkan
  $newName = uniqid('file_') . '.' . $ext;
  if (!is_dir($uploadDir)) mkdir($uploadDir, 0755, true);

  if (move_uploaded_file($file['tmp_name'], $uploadDir . $newName)) {
    echo "Upload berhasil: {$newName}";
  } else {
    echo "Gagal memindahkan file.";
  }
}
?>
\end{webcode}

